\section{Introduction}
Multi-access edge computing (MEC) places compute close to data producers while
retaining access to larger cloud or HPC resources \cite{etsiMec003}. A task
must therefore be mapped to a device CPU, an associated edge CPU, or a shared
cloud queue. The choice couples execution time, queueing, wireless and
backhaul transmission, device energy, and a relative deadline. A centralized
controller can observe more state, but it also requires broader information
sharing and can become a coordination bottleneck.

This paper positions the repository as an experimental framework for studying
these choices under common simulator conditions. Each tier makes only the
decision exposed by its local contract: the device chooses local execution or
its assigned edge, the edge chooses edge execution or cloud forwarding, and
the cloud selects a queue discipline. A common SimPy backend executes the
centralized and hierarchical policies with shared task, timing, energy, and
metric definitions, making cross-policy differences directly inspectable.

The paper makes five contributions. First, it implements a three-tier
device--edge--cloud architecture with distinct decision boundaries and a
shared cloud scheduler. Second, it integrates centralized DQN, hierarchical
DQN, and hierarchical Double DQN controllers into that common backend.
Third, it implements causal adjacent-tier feedback with modeled bytes, delay,
and information age. Fourth, it reports a matched comparison on 579
chronologically held-out Alibaba-derived tasks using 24 devices, three edges,
and one cloud worker. Fifth, it analyzes deadline success, latency, device
energy, communication overhead, feedback strategies, topology changes, and
cloud scheduling. The results expose objective-dependent trade-offs rather
than assuming that one controller dominates every metric. The learning
formulation and workload adaptation define the scope of the evidence; they
also provide a controlled basis for future sequential multi-agent studies.
